Przedstawimy teraz algorytm znajdujący najliczniejsze skojarzenie w grafie dwudzielnym (Problem~\ref{problem:matching}).

Niech $M$ będzie skojarzeniem w $G$. \\
Definiujemy graf skierowany $G_M = (V_1 \cup V_2, E_M)$, gdzie
\begin{flalign*}
	E_M = \{(u, v) : uv \in E - M,\ u \in V_1,\ v \in V_2\} \cup \{(v, u) : uv \in M,\ u \in V_1,\ v \in V_2\}. &  &
\end{flalign*}
Innymi słowy, krawędzie wolne kierujemy z $V_1$ do $V_2$, a krawędzie skojarzone z $V_2$ do $V_1$.

\begin{example} ~\\
	\begin{minipage}{0.48\textwidth}
		\centering
		\begin{tikzpicture}[x=0.9cm, y=0.9cm]
			\foreach \x/\name in {0/a, 1/b, 2/c, 3/d, 4/e} {
					\node[circle, fill=black, inner sep=1.5pt] (u\name) at (\x, 1) {};
				}
			\foreach \x/\name in {0/p, 1/q, 2/r, 3/s, 4/t} {
					\node[circle, fill=black, inner sep=1.5pt] (l\name) at (\x, 0) {};
				}
			\node[right] at (4.3, 1) {$V_1$};
			\node[right] at (4.3, 0) {$V_2$};
			\node[left] at (-0.3, 0.5) {$G$};


			\draw (ua) -- (lq);
			\draw (ub) -- (lp);
			\draw (uc) -- (lp);
			\draw[line width=1.5pt, accent] (uc) -- (lq);
			\draw (uc) -- (lr);
			\draw (uc) -- (ls);
			\draw (ud) -- (lr);
			\draw[line width=1.5pt, accent] (ud) -- (lt);
			\draw (ue) -- (ls);
			\draw (ue) -- (lt);
		\end{tikzpicture}
	\end{minipage}\hfill
	\begin{minipage}{0.48\textwidth}
		\centering
		\begin{tikzpicture}[x=0.9cm, y=0.9cm, >=stealth]
			\foreach \x/\name in {0/a, 1/b, 2/c, 3/d, 4/e} {
					\node[circle, fill=black, inner sep=1.5pt] (u\name) at (\x, 1) {};
				}
			\foreach \x/\name in {0/p, 1/q, 2/r, 3/s, 4/t} {
					\node[circle, fill=black, inner sep=1.5pt] (l\name) at (\x, 0) {};
				}
			\node[right] at (4.3, 1) {$V_1$};
			\node[right] at (4.3, 0) {$V_2$};
			\node[left] at (-0.3, 0.5) {$G_M$};

			\draw[->] (ua) -- (lq);
			\draw[->] (ub) -- (lp);
			\draw[->] (uc) -- (lp);
			\draw[<-, line width=1.2pt, accent] (uc) -- (lq);
			\draw[->] (uc) -- (lr);
			\draw[->] (uc) -- (ls);
			\draw[->] (ud) -- (lr);
			\draw[<-, line width=1.2pt, accent] (ud) -- (lt);
			\draw[->] (ue) -- (ls);
			\draw[->] (ue) -- (lt);
		\end{tikzpicture}
	\end{minipage}
\end{example}

Procedura do znajdowania ścieżki powiększającej --- \textit{Augmenting Path Search} (APS).

\begin{alg}{Augmenting Path Search}{alg:aps:lect14}
	\FunctionNN{APS}{$G, M$} \Comment{$G = (V_1, V_2; E)$, $M$ -- skojarzenie w $G$}
	\State $V_1'$ \Assign zbiór wierzchołków wolnych z $V_1$
	\State $V_2'$ \Assign zbiór wierzchołków wolnych z $V_2$
	\State skonstruuj graf skierowany $G_M = (V_1 \cup V_2, E_M)$
	\State znajdź ścieżkę $\vec{P}$ z $V_1'$ do $V_2'$ w $G_M$
	\If{$\vec{P}$ istnieje}
	\State \Return $P$ \Comment{$P$ -- nieskierowana $\vec{P}$}
	\Else
	\State \Return NULL
	\EndIf
	\EndFunction
\end{alg}

\begin{lemma}
	Procedura \textnormal{\Call{APS}{}} zwraca ścieżkę $P \iff$ istnieje ścieżka powiększająca względem $M$ w $G$. Ponadto ścieżka $P$ zwrócona przez \textnormal{\Call{APS}{}} jest ścieżką powiększającą względem $M$.
\end{lemma}

\begin{proof} ~
	\begin{itemize}
		\item[$\implies$] Jeśli \textnormal{\Call{APS}{}} zwraca ścieżkę $P$, to skierowana ścieżka $\vec{P}$ znaleziona w linii 4 zaczyna się w wierzchołku z $V_1'$ (a więc wolnym) i kończy się w wierzchołku z $V_2'$ (a więc wolnym). Z definicji grafu $G_M$ wynika, że pierwsza krawędź $P$ jest wolna, druga krawędź $P$ jest skojarzona, trzecia krawędź $P$ jest wolna, \dots, ostatnia krawędź $P$ jest wolna. Zatem $P$ jest ścieżką powiększającą względem $M$.

		\item[$\impliedby$] Jeśli istnieje ścieżka powiększająca względem $M$ w $G$, to ścieżka skierowana otrzymana z takiej ścieżki powiększającej przez skierowanie krawędzi zgodnie z $E_M$ jest ścieżką skierowaną z $V_1'$ do $V_2'$ w $G_M$. Taką ścieżkę skierowaną znajduje \textnormal{\Call{APS}{}}.
	\end{itemize}
\end{proof}

Algorytm znajdowania najliczniejszego skojarzenia --- \textit{Maximum Matching} (MM).

\begin{alg}{Maximum Matching}{alg:mm:lect14}
	\FunctionNN{MM}{$G$} \Comment{$G = (V_1, V_2; E)$}
	\State $M \Assign \emptyset$
	\Repeat
	\State $P \Assign \Call{APS}{G, M}$
	\If{$P \neq$ NULL}
	\State $M \Assign M \div E(P)$
	\EndIf
	\Until{$P =$ NULL}
	\State \Return $M$
	\EndFunction
\end{alg}

\paragraph{Poprawność} Wynika z twierdzenia Berge'a.

\paragraph{Złożoność czasowa}

Złożoność czasowa \textnormal{\Call{APS}{}}:
\begin{itemize}[nosep]
	\item Linie 1, 2: $\BO(|V|)$.
	\item Linia 3: $\BO(|E|)$.
	\item Linia 4: $\BO(|E|)$ (używamy algorytmu BFS).
	\item Linie 5--8: $\BO(|V|)$.
\end{itemize}
Złożoność \textnormal{\Call{APS}{}} wynosi zatem $\BO(|E|)$.

Złożoność czasowa \textnormal{\Call{MM}{}}:
\begin{itemize}[nosep]
	\item Linia 1: $\BO(1)$.
	\item Pętla 2--6 wykona się co najwyżej $\frac{|V|}{2}$ razy, bo w każdej iteracji $|M|$ zwiększa się o $1$, a maksymalna liczność skojarzenia jest $\le \frac{|V|}{2}$.
	\item Wnętrze pętli:
	      \begin{itemize}[nosep]
		      \item Linia 3: $\BO(|E|)$.
		      \item Linia 5: $\BO(|V|)$.
	      \end{itemize}
	      Razem $\BO(|E|)$.
	\item Linia 7: $\BO(|V|)$.
\end{itemize}
Ostateczna złożoność \textnormal{\Call{MM}{}} wynosi $\BO(|V| \cdot |E|)$.

% Wspólne makra rysunkowe dla przykładu MM.
% Graf G (kolejność krawędzi): aq, bp, cp, cq, cr, cs, dr, dt, es, et.
\tikzset{
	mmN/.style={circle, fill=black, inner sep=1.5pt},
	matN/.style={<-},                          % skojarzona: strzałka w górę (V_2->V_1)
	augFN/.style={path, line width=1.2pt},     % na ścieżce, wolna (w dół, domyślne ->)
	augMN/.style={path, line width=1.2pt, <-}, % na ścieżce, skojarzona (w górę)
	pthN/.style={path, line width=1.2pt},      % wyróżnienie w nieskierowanym G
	o/.style={},                               % pusty styl (pgffor gubi pusty pierwszy człon)
}
\newcommand{\mmnodes}{%
	\foreach \x/\n in {0/a, 1/b, 2/c, 3/d, 4/e} {\node[mmN] (u\n) at (\x, 1) {};}
	\foreach \x/\n in {0/p, 1/q, 2/r, 3/s, 4/t} {\node[mmN] (l\n) at (\x, 0) {};}
	\node[right] at (4.3, 1) {$V_1$}; \node[right] at (4.3, 0) {$V_2$};}
\newcommand{\mmfree}[1]{\ifx\relax#1\relax\else
	\foreach \n in {#1} {\draw[draw=black] (\n) circle (4pt);}\fi}
% Krawędzie. #1: lista par "styl/od/do" (TeX dopuszcza max 9 parametrów,
% więc krawędzie idą przez \foreach zamiast osobnych argumentów).
\newcommand{\mmarcs}[1]{\foreach \st/\a/\b in {#1} {\draw[->,\st] (\a)--(\b);}}
\newcommand{\mmlines}[1]{\foreach \st/\a/\b in {#1} {\draw[\st] (\a)--(\b);}}
% Panel G_M: #1 lista par styl/od/do, #2 lista wolnych, #3 wierzchołek startowy (różowy pierścień, pusty dozwolony).
\newcommand{\mmGM}[3]{%
	\begin{tikzpicture}[x=0.9cm, y=0.9cm, >=stealth, baseline=(current bounding box.center)]
		\mmnodes \node[left] at (-0.3, 0.5) {$G_M$};
		\mmfree{#2}%
		\ifx\relax#3\relax\else\draw[path, line width=1.2pt] (#3) circle (4pt);\fi
		\mmarcs{#1}\end{tikzpicture}}
% Panel G z wyróżnionym skojarzeniem: #1 lista par styl/od/do.
\newcommand{\mmG}[1]{%
	\begin{tikzpicture}[x=0.9cm, y=0.9cm, >=stealth, baseline=(current bounding box.center)]
		\mmnodes \node[left] at (-0.3, 0.5) {$G$};
		\mmlines{#1}\end{tikzpicture}}

\begin{example}[Działanie algorytmu Maximum Matching]
	Prześledźmy działanie algorytmu \textnormal{\Call{MM}{}} na poniższym grafie $G$.
	W~każdej iteracji po~lewej $G_M$ ze~ścieżką powiększającą
	(\textcolor{path}{na~różowo}, różowy pierścień to~wierzchołek startowy),
	po~prawej graf $G$ z~bieżącym skojarzeniem $M$ (\textcolor{path}{wyróżnionym}).

	\begin{azfigure}
		\centering
		\mmG{o/ua/lq, o/ub/lp, o/uc/lp, o/uc/lq, o/uc/lr, o/uc/ls, o/ud/lr, o/ud/lt, o/ue/ls, o/ue/lt}
	\end{azfigure}

	\textbf{Iteracja 1:}
	\begin{azfigure}
		\centering
		\mmGM{o/ua/lq, o/ub/lp, o/uc/lp, augFN/uc/lq, o/uc/lr, o/uc/ls, o/ud/lr, o/ud/lt, o/ue/ls, o/ue/lt}{ua,ub,uc,ud,ue,lp,lq,lr,ls,lt}{uc}\qquad
		\mmG{o/ua/lq, o/ub/lp, o/uc/lp, pthN/uc/lq, o/uc/lr, o/uc/ls, o/ud/lr, o/ud/lt, o/ue/ls, o/ue/lt}
	\end{azfigure}

	\textbf{Iteracja 2:}
	\begin{azfigure}
		\centering
		\mmGM{o/ua/lq, o/ub/lp, o/uc/lp, matN/uc/lq, o/uc/lr, o/uc/ls, o/ud/lr, o/ud/lt, augFN/ue/ls, o/ue/lt}{ua,ub,ud,ue,lp,lr,ls,lt}{ue}\qquad
		\mmG{o/ua/lq, o/ub/lp, o/uc/lp, pthN/uc/lq, o/uc/lr, o/uc/ls, o/ud/lr, o/ud/lt, pthN/ue/ls, o/ue/lt}
	\end{azfigure}

	\textbf{Iteracja 3:}
	\begin{azfigure}
		\centering
		\mmGM{o/ua/lq, o/ub/lp, o/uc/lp, matN/uc/lq, o/uc/lr, o/uc/ls, o/ud/lr, augFN/ud/lt, matN/ue/ls, o/ue/lt}{ua,ub,ud,lp,lr,lt}{ud}\qquad
		\mmG{o/ua/lq, o/ub/lp, o/uc/lp, pthN/uc/lq, o/uc/lr, o/uc/ls, o/ud/lr, pthN/ud/lt, pthN/ue/ls, o/ue/lt}
	\end{azfigure}

	\textbf{Iteracja 4:}
	\begin{azfigure}
		\centering
		\mmGM{augFN/ua/lq, o/ub/lp, o/uc/lp, augMN/uc/lq, o/uc/lr, augFN/uc/ls, augFN/ud/lr, augMN/ud/lt, augMN/ue/ls, augFN/ue/lt}{ua,ub,ud,lp,lr,lt}{ua}\qquad
		\mmG{pthN/ua/lq, o/ub/lp, o/uc/lp, o/uc/lq, o/uc/lr, pthN/uc/ls, pthN/ud/lr, o/ud/lt, o/ue/ls, pthN/ue/lt}
	\end{azfigure}

	\textbf{Iteracja 5:}
	\begin{azfigure}
		\centering
		\mmGM{matN/ua/lq, augFN/ub/lp, o/uc/lp, o/uc/lq, o/uc/lr, matN/uc/ls, matN/ud/lr, o/ud/lt, o/ue/ls, matN/ue/lt}{ub,lp}{ub}\qquad
		\mmG{pthN/ua/lq, pthN/ub/lp, o/uc/lp, o/uc/lq, o/uc/lr, pthN/uc/ls, pthN/ud/lr, o/ud/lt, o/ue/ls, pthN/ue/lt}
	\end{azfigure}
\end{example}
